\section{Method}
\label{sec:method}

Our pipeline transforms a single RGB input image into a Kusama-style
polka dot rendering through five stages: (1) monocular depth
estimation, (2) semantic segmentation, (3) depth-density field
construction, (4) three-layer depth-aware stippling, and (5)
semantic-aware rendering.

\subsection{Monocular Depth Estimation}
\label{sec:method:depth}

We employ the YOLO26x-depth model, the largest scale (x) of the YOLO26
depth estimation family, to obtain a dense per-pixel metric depth map.
Given an input image $I \in \mathbb{R}^{H \times W \times 3}$, the
model produces a depth map $D \in \mathbb{R}^{H' \times W'}$ where each
value represents the estimated distance to the scene point in meters.
We resize $D$ to the original image resolution via bicubic interpolation
to obtain $D \in \mathbb{R}^{H \times W}$.

\subsubsection{Depth Normalization}

Raw metric depth values can span a wide range, with outlier pixels
(e.g., sky at near-infinite distance) compressing the useful depth
range. A naive global normalization over the entire image would map
most of the $[0, 1]$ range to background depth variation (sky, distant
buildings), leaving only a narrow band for the foreground objects that
are the artistic focus. To ensure that the full normalized depth range
captures depth variation \emph{within} foreground objects, we employ
\emph{foreground-aware depth normalization}:

\begin{equation}
    \hat{d}(x,y) = \begin{cases}
        \displaystyle\frac{D - d_{\min}^{\text{fg}}}{d_{\max}^{\text{fg}} - d_{\min}^{\text{fg}}} & \text{if } M_{\text{fg}} = 1 \\[6pt]
        1.0 & \text{otherwise}
    \end{cases}
    \label{eq:depth_norm}
\end{equation}

\noindent where $D \equiv D(x,y)$, $d_{\min}^{\text{fg}}$ and $d_{\max}^{\text{fg}}$ are
the $p$-th and $(100{-}p)$-th percentiles of the depth values within
the foreground mask $M_{\text{fg}}$ (defined in
Section~\ref{sec:method:seg}), with $p = 2$ by default. The result is
clipped to $[0, 1]$. Background
pixels are assigned $\hat{d} = 1.0$ (farthest), so they receive the
smallest dots.

This design amplifies subtle depth differences within foreground
objects---e.g., the front versus the back of a pumpkin, or the nose
versus the ears of a cat---into the full $[0, 1]$ normalized range,
producing visible dot size variation that reflects the object's
three-dimensional form. In our experiments, foreground-aware
normalization increases the depth standard deviation within
foreground regions by approximately 45\% compared to global
normalization (from $\sigma_{\hat{d}} = 0.158$ to $0.229$), directly
translating to richer dot size diversity on the main subject.

\subsection{Semantic Segmentation}
\label{sec:method:seg}

We employ the YOLO26x-sem model to obtain a per-pixel semantic
segmentation map $S \in \{0, 1, \ldots, 18\}^{H \times W}$, where each
value is a Cityscapes class ID. We define a set of background classes
$\mathcal{B} = \{\textit{road}, \textit{sidewalk}, \textit{building},
\textit{wall}, \textit{sky}, \textit{vegetation}, \textit{terrain}\}$
and compute a binary foreground mask:

\begin{equation}
    M_{\text{fg}}(x,y) = \begin{cases} 1 & \text{if } S(x,y) \notin \mathcal{B} \\ 0 & \text{otherwise} \end{cases}
    \label{eq:fg_mask}
\end{equation}

Small connected components in $M_{\text{fg}}$ (area $< 2\%$ of canvas)
are removed to eliminate noise fragments. This foreground mask enables
dual-color rendering (Section~\ref{sec:method:rendering}).

\subsection{Depth-Density Field Construction}
\label{sec:method:density}

The density field $\rho(x,y)$ drives Voronoi stippling point placement:
regions with higher density receive more dots. We construct a composite
density field that combines depth and luminance information:

\begin{equation}
    \rho(x,y) = \underbrace{(1 - \hat{d}(x,y))^\beta}_{\text{depth density}} \cdot \underbrace{\gamma(I_{\text{gray}}(x,y))}_{\text{luminance density}},
    \label{eq:composite_density}
\end{equation}

\noindent where $\beta$ is the depth-density exponent (default 0.6),
$\gamma(\cdot)$ is the gamma-corrected inverted luminance density from
standard stippling \citep{secord2002weighted}, and $I_{\text{gray}}$ is
the grayscale conversion of the input image. The depth density term
$(1 - \hat{d})^\beta$ concentrates dots in near regions, while the
luminance density $\gamma$ preserves the traditional dark-region
emphasis. The exponent $\beta < 1$ ensures that far regions remain
populated with sparse dots rather than being entirely empty, preserving
the horror-vacui aesthetic.

The composite density is normalized to $[0, 1]$ before use.

\subsection{Depth-to-Radius Mapping}
\label{sec:method:radius}

The central mathematical component of this work is the
depth-to-radius mapping function that assigns each dot a radius based
on the local depth value:

\begin{equation}
    r(x,y) = \Bigl[\, r_{\min} + (r_{\max} - r_{\min}) \cdot (1 - \hat{d})^\alpha \,\Bigr] \cdot (1 + \varepsilon),
    \label{eq:radius_mapping}
\end{equation}

\noindent where:
\begin{itemize}
    \item $r_{\min}$ and $r_{\max}$ are the minimum and maximum dot
    radii (in pixels), defining the multi-scale range.
    \item $\alpha > 0$ is the \emph{depth falloff exponent} controlling
    the curvature of the depth-size relationship. When $\alpha < 1$,
    the mapping is concave, producing gentle size variation where most
    dots are medium-sized. When $\alpha = 1$, the mapping is linear.
    When $\alpha > 1$, the mapping is convex, producing dramatic size
    variation with strong depth emphasis---large dots in near regions,
    rapidly diminishing to small dots in mid-depth.
    \item $\varepsilon \sim \mathcal{N}(0, \sigma^2)$ is a
    \emph{multiplicative} Gaussian jitter with standard deviation
    $\sigma$ (default 0.15) that injects organic variation. Unlike an
    additive noise term, the multiplicative form scales the absolute
    jitter with the nominal radius, so small dots remain small and
    large dots exhibit proportionally larger irregularities---matching
    the statistics of hand-painted marks. The clipped output keeps
    radii within $[\,0.3\,\mathrm{px},\; 1.1\, r_{\max}\,]$.
\end{itemize}

This formulation has several desirable properties:
\begin{proposition}[Monotonicity]
The expected dot radius $\mathbb{E}[r]$ is a monotonically decreasing
function of depth $\hat{d}$ for all $\alpha > 0$, ensuring that near
objects always receive larger dots on average.
\end{proposition}

\begin{proposition}[Boundedness]
The dot radius is bounded with high probability:
$\Pr\!\bigl[\, r \in [\,r_{\min}(1 - 3\sigma),\; r_{\max}(1 + 3\sigma)\,] \,\bigr] > 99.7\%$,
so no dot exceeds a predictable maximum size.
\end{proposition}

The parameter $\alpha$ provides intuitive artistic control: higher
values produce more dramatic depth emphasis (a few very large dots in
near regions, many tiny dots elsewhere), while lower values produce a
more uniform, gentle gradient of sizes. The noise parameter $\sigma$
controls the degree of organic variation: $\sigma = 0$ yields perfectly
deterministic size assignment (mechanical appearance), while $\sigma >
0.2$ produces high variability (chaotic, expressionistic appearance).

\subsection{Learnable Depth-to-Radius Mapping}
\label{sec:method:learned}

While the fixed power-law mapping (Eq.~\ref{eq:radius_mapping}) provides
interpretable control, the optimal depth-to-size relationship may vary
across scene types and artistic preferences. We introduce a
\emph{learnable} alternative: a lightweight CNN
(${\sim}47$\,K parameters, 4 conv layers) that takes per-pixel
features---normalized depth $\hat{d}$, luminance $I_{\text{gray}}$,
depth-gradient magnitude $|\nabla\hat{d}|$, and luminance-gradient
magnitude $|\nabla I|$---and predicts both a radius map and a density
weight map.

The network is trained via a \emph{perceptual stippling loss}
combining four terms: (i)~depth-gradient alignment (radius should
correlate with $(1-\hat{d})$), (ii)~coverage matching (target
${\sim}28\%$), (iii)~radius-histogram matching (target distribution
skewed toward small dots with a long tail, matching Kusama's
statistics), and (iv)~spatial smoothness (total-variation
regularization). Training is self-supervised and requires no
artist-annotated ground truth. The learned mapping can be
seamlessly swapped into the pipeline as a drop-in replacement for
Eq.~\ref{eq:radius_mapping}, with the same multiplicative Gaussian
jitter for organic variation.

This component addresses a limitation of purely algorithmic NPR
methods: by learning the depth-size relationship from data, the
pipeline adapts to diverse scene types without manual parameter
tuning, providing the data-driven contribution that distinguishes
our approach from classical stippling.

\subsection{Three-Layer Stippling Architecture}
\label{sec:method:stippling}

A fundamental tension exists in depth-aware stippling: large dots
(required for near-region emphasis) occupy significant canvas area and
can merge with nearby dots if not adequately spaced. Single-layer
approaches that assign a wide radius range ($r_{\min} = 1$\,px to
$r_{\max} = 25$\,px) to all points suffer from blob formation, as the
overlap-pruning step must choose between keeping large dots (which
dominate the composition) or keeping small dots (which provide texture).

We resolve this tension through a three-layer architecture, each layer
operating in a distinct radius regime with independent overlap
pruning:

\subsubsection{Layer 1: Fine Stipple}

The fine stipple layer provides the base texture: a dense field of
small dots ($r \in [1, 5]$\,px) whose density follows the composite
depth-luminance field (Eq.~\ref{eq:composite_density}). Points are
initialized by rejection sampling, relaxed via Lloyd iteration
(8 iterations with jitter to break hexagonal regularity), and
assigned radii via the depth-to-radius mapping
(Eq.~\ref{eq:radius_mapping}) with a reduced falloff exponent
$\alpha' = 0.6\alpha$ for subtle size variation. Overlap pruning
uses a tight gap of 1.0\,px, allowing dense packing of small dots.

\subsubsection{Layer 2: Depth-Emphasis}

The depth-emphasis layer carries the primary depth signal: large dots
($r \in [7.5, r_{\max}]$\,px) placed exclusively in near foreground
regions, gated by the condition
$M_{\text{near}} = \mathbb{1}[\hat{d} < 0.55] \cdot M_{\text{fg}}$.
Points are sampled from a near-weighted density field
$\rho_{\text{near}} = M_{\text{near}} \cdot \rho$. Overlap pruning
uses a larger gap of 4.0\,px to ensure adequate spacing between large
dots.

\paragraph{Cross-Layer Pruning.}
After generating both layers, we perform cross-layer pruning: fine
dots whose centers fall within the bounding circle of any
depth-emphasis dot (plus a 1\,px margin) are removed. This prevents
large dots from connecting to nearby fine dots and forming merged
blobs. The pruning is implemented via a $k$-d tree for efficient
neighbor queries:

\begin{algorithm}
\caption{Cross-Layer Overlap Pruning}
\label{alg:cross_prune}
\begin{algorithmic}[1]
\REQUIRE Fine dots $\{(p_i, r_i)\}$, emphasis dots $\{(q_j, s_j)\}$
\ENSURE Pruned fine dots $\{(p_i, r_i)\}$
\STATE Build $k$-d tree $T$ on $\{q_j\}$
\FOR{each fine dot $(p_i, r_i)$}
    \STATE $R \leftarrow r_i + \max_j s_j + 2$ \COMMENT{search radius}
    \STATE $\mathcal{N} \leftarrow T.\text{query}(p_i, R)$
    \STATE $\text{keep}_i \leftarrow \text{true}$
    \FOR{each $j \in \mathcal{N}$}
        \IF{$\|p_i - q_j\| < r_i + s_j + 1$}
            \STATE $\text{keep}_i \leftarrow \text{false}$;\; \textbf{break}
        \ENDIF
    \ENDFOR
\ENDFOR
\RETURN $\{(p_i, r_i) : \text{keep}_i = \text{true}\}$
\end{algorithmic}
\end{algorithm}

\subsubsection{Layer 3: Edge and Background}

The edge layer places small dots ($r = 1.3$\,px) at regular intervals
along Canny edges \citep{canny1986computational} with Gaussian
jitter, providing structural emphasis along object contours. The
background layer applies a sparse jittered grid of tiny dots ($r \approx
1$\,px) in low-density regions (composite density $< 0.25$), satisfying
the horror-vacui requirement that no canvas region remains empty.

\subsection{Semantic-Aware Rendering}
\label{sec:method:rendering}

The rendering stage draws all dots on a flat background canvas using
hard-edged solid circles (\texttt{cv2.circle} with
\texttt{LINE\_4} connectivity, no anti-aliasing). The hard-edge
property is both stylistically correct---Kusama's dots have crisp,
mechanical boundaries---and functionally necessary, as anti-aliasing
would create intermediate-alpha bridge pixels between adjacent dots,
causing them to visually merge.

Each Kusama palette defines three colors: a background color
$c_{\text{bg}}$, a foreground dot color $c_{\text{fg}}$, and a
background dot color $c_{\text{bg-dot}}$. When semantic-aware mode is
enabled, each dot's color is determined by its center pixel in the
foreground mask:

\begin{equation}
    c_i = \begin{cases} c_{\text{fg}} & \text{if } M_{\text{fg}}(y_i, x_i) = 1 \\ c_{\text{bg-dot}} & \text{otherwise} \end{cases}
    \label{eq:semantic_color}
\end{equation}

This creates a figure--ground color separation: foreground objects
are rendered in one dot color while
background regions use a contrasting dot color,
both against a shared flat background. The five palettes are designed
with Kusama's characteristic high chroma and flat application:

\begin{itemize}
    \item \textbf{black\_red}: black background, red foreground dots,
    cyan background dots (classic Kusama)
    \item \textbf{yellow\_black}: chrome yellow background, black
    foreground dots, dark red background dots (pumpkin series)
    \item \textbf{white\_red}: white background, scarlet foreground
    dots, cyan-blue background dots (high contrast)
    \item \textbf{magenta\_cyan}: near-black background, magenta
    foreground dots, cyan background dots
    \item \textbf{red\_yellow}: red background, yellow foreground dots,
    white background dots (inverse configuration)
\end{itemize}

The drawing order ensures visual hierarchy: background fill dots are
drawn first (underneath), followed by the fine stipple layer, edge
dots, and finally the depth-emphasis layer on top (most prominent).

\subsection{Implementation Details}

The pipeline is implemented in Python using PyTorch for YOLO26
inference, OpenCV for image processing, and SciPy for Voronoi
computation. The YOLO26x-depth and YOLO26x-sem models run on an NVIDIA
RTX A4500 GPU (20\,GB VRAM), with inference times of approximately
2.2\,s and 0.4\,s respectively for a $1024 \times 1024$ image. The
Voronoi stippling stage (rejection sampling, Lloyd relaxation, and
overlap pruning) runs on CPU and takes approximately 60\,s for 25,000
initial points with 8 relaxation iterations. The Lloyd relaxation uses
the mirror-points trick to handle unbounded Voronoi cells and
cumulative-sum acceleration for density-weighted centroid computation,
following \citet{secord2002weighted}.
